\documentclass[11pt]{article}
\usepackage[a4paper,margin=25mm]{geometry}
\usepackage[T1]{fontenc}
\usepackage{lmodern,amsmath,amssymb,amsthm,mathtools,booktabs,microtype}
\usepackage{enumitem}
\usepackage{titlesec}
\titleformat{\part}[display]{\Large\bfseries}{Part \thepart}{0.6ex}{}
\titlespacing*{\part}{0pt}{4ex plus 1ex}{2ex}
\usepackage[noend]{algpseudocode}
\usepackage[dvipsnames]{xcolor}
\usepackage[colorlinks=true,linkcolor=MidnightBlue,urlcolor=MidnightBlue,citecolor=MidnightBlue]{hyperref}
\hypersetup{pdftitle={A Boolean linear-time planarity test},pdfauthor={MohammadTaghi Hajiaghayi and Mahdi JafariRaviz}}
\setlength{\parskip}{4pt}
\setlength{\parindent}{0pt}
\setlength{\emergencystretch}{2em}

% Sections are numbered from 1 inside every part. Subsections of Part IV,
% Section 7 are numbered 7a, 7b and 7c.
\counterwithin*{section}{part}
\renewcommand{\thesubsection}{\arabic{section}\alph{subsection}}
\renewcommand{\theHsection}{\thepart.\arabic{section}}
\renewcommand{\theHsubsection}{\thepart.\arabic{section}.\arabic{subsection}}

\newtheorem*{theorem}{Theorem}
\newcommand{\lemmaname}{}
\newtheorem*{namedlemma}{Lemma \lemmaname}
\newenvironment{lem}[1]{\renewcommand{\lemmaname}{#1}\begin{namedlemma}}{\end{namedlemma}}

\DeclareMathOperator{\low}{low}
\DeclareMathOperator{\lowtwo}{low2}
\DeclareMathOperator{\high}{high}
\DeclareMathOperator{\thick}{thick}
\DeclareMathOperator{\chordal}{chordal}
\DeclareMathOperator{\key}{key}
\DeclareMathOperator{\subtree}{subtree}
\DeclareMathOperator{\pos}{pos}
\DeclareMathOperator{\maxkey}{max}
\DeclareMathOperator{\minkey}{min}
\newcommand{\cL}{\mathsf{L}}
\newcommand{\cR}{\mathsf{R}}
\newcommand{\pr}[1]{\textsc{#1}}

\title{\textbf{A Boolean linear-time planarity test}}
\author{MohammadTaghi Hajiaghayi \and Mahdi JafariRaviz}
\date{November 2026}

\begin{document}
\maketitle

\begin{abstract}
We give a deterministic program that decides in time $O(N+M+1)$ whether a
finite simple graph with $N$ vertices and $M$ edges is planar. It returns only
the answer, not an embedding. The program is a left-right planarity test in the
tradition of de Fraysseix and Rosenstiehl. The note proves everything it uses:
the correctness of the constraint solver, and the left-right characterization
of planarity itself. The only fact about planar graphs that we use is Euler's
edge bound. The test supplies the planarity primitive of a companion note on
extracting crossings, which replaces the two paths algorithm of Shiloach in
Graph Minors XIII.
\end{abstract}

\section*{Introduction}

All graphs in this note are finite and simple. A \emph{plane embedding} of a
graph places its vertices at distinct points of the plane and draws each edge
as an injective continuous arc between the points of its two ends. Two edge
images meet only at common ends, and no edge image contains a vertex point
other than its own ends. A graph is \emph{planar} if it has a plane embedding.
Part~\ref{part:iv}, Section~6 also treats embeddings given combinatorially, by
rotation systems.

\begin{theorem}
There is a single deterministic word-RAM program which, given the adjacency
lists of a graph with $N$ vertices and $M$ edges, returns true exactly when the
graph has a plane embedding. Its time and space are $O(N+M+1)$. The program
need not return an embedding.
\end{theorem}

\paragraph{How the test works.}
Run a depth-first search. In a plane embedding, every back edge leaves the
search tree either on the left or on the right of the tree path it closes.
De Fraysseix and Rosenstiehl showed that a graph is planar if and only if its
back edges can be split into left and right ones so that certain sets of back
edges get one side and certain pairs of sets get opposite sides
\cite{dFR82,dFR85}. In the form of de Fraysseix, Ossona de Mendez and
Rosenstiehl \cite{dFOMR06} (Theorem~2 on page~8 of the arXiv version), and in
Brandes's description of the left-right planarity test \cite[Theorem~5,
page~10]{Br09}, these are constraints on two colors $\cL$ and $\cR$ of the back
edges. Part~\ref{part:i} of this note decides these constraints in one pass
over the search tree. Each vertex hands its parent a compact normal form of
the constraints found below it: a stack of blocks, each block made of two
lists of back edges. We add one extra condition, weak alignment, which keeps
this normal form simple. It does not change whether the constraints can be
satisfied (Part~\ref{part:iii}). Brandes uses a related alignment of the
partition \cite[Lemma~6, page~11]{Br09}.

The program uses depth-first search, integer bucket sorting and linked lists.
No planarity test is assumed. We prove every characterization and every
property of the data structure that we use. The cited works are background,
not premises of the theorem. The only fact about planar graphs that we use is
Euler's bound: a simple planar graph with $N\ge3$ vertices has at most $3N-6$
edges. It is used in Part~\ref{part:iv}, Section~\ref{sec:k5}. Part~\ref{part:iv},
Section~6 gives a second proof of the same step for rotation systems, which
uses Euler's formula instead.

\paragraph{Where the test is used.}
The companion note \cite{HJ26} gives an algorithm that extracts a crossing of a
cycle in time $O(nm)$. It calls a Boolean linear-time planarity test, which it
states as Fact~P. That algorithm replaces Shiloach's two paths algorithm
\cite{Sh80} in (8.1) of Graph Minors XIII \cite[page~93]{GM13}. Linear-time
planarity testing goes back to Hopcroft and Tarjan \cite{HT74}. In the Lean
formalization of the Graph Minors series, the test of this note proves
Fact~P, so the main result of Graph Minors XIII is proved there without
assuming a planarity algorithm (Part~\ref{part:vi}).

\paragraph{Outline.}
Part~\ref{part:i} gives the algorithm. Part~\ref{part:ii} proves that it
decides exactly the constraints stated in Part~\ref{part:i}, Section~2.
Part~\ref{part:iii} shows that weak alignment does not change satisfiability.
Part~\ref{part:iv} proves that the constraints can be satisfied exactly when
the graph is planar. Part~\ref{part:v} bounds the running time on a word RAM,
assembles the correctness proof, explains the use in the crossing extractor,
and describes tests of the reference implementation
\texttt{planarity\_linear.py}. Part~\ref{part:vi} describes the Lean
formalization.

%======================================================================
\part{The exact algorithm}\label{part:i}

This part gives the whole algorithm. Section~1 runs a depth-first search and
computes for every edge a key, which fixes the order in which edges are
processed. Section~2 states the constraints on two-colorings of the back edges
that the algorithm decides. Sections~3 to~6 give the data structure and the
procedure that decides them.

\section{DFS definitions and preprocessing}

We run an ordinary depth-first search of each component. We orient tree edges
from parent to child and all other edges from descendant to ancestor. Let
$h(v)$ be the depth of $v$, with root depth zero, and let $\subtree(v)$ be the
set of descendants of $v$, including $v$. All non-tree edges have comparable
ends: if an undirected edge joined two different branches of the search, then
scanning it from the first branch would discover the other end before that
branch is finished, a contradiction. We call the non-tree edges \emph{back
edges}.

For a back edge $b=(x,u)$, put $d(b)=h(u)$. The \emph{outgoing edges} of a
vertex $v$ are the tree edges from $v$ to its children and the back edges
from $v$. For an outgoing edge $e$ of $v$, define
\begin{align*}
F(e)&=\{e\}, &&\text{if $e$ is a back edge,}\\
F(e)&=\{\text{back edges }(x,u):\ x\in\subtree(w),\ h(u)<h(v)\},
&&\text{if $e=(v,w)$ is a tree edge.}
\end{align*}
These are the \emph{strict returns} of $e$. Define
\begin{align*}
\low(e)&=\min\bigl(\{h(v)\}\cup\{d(b):b\in F(e)\}\bigr),\\
\thick(e)&=1\text{ if $F(e)$ has two different $d$-values, and $0$ otherwise,}\\
\key(e)&=2\low(e)+\thick(e).
\end{align*}
We also call a thick edge \emph{chordal}, and write
$\chordal(e)=\thick(e)$. We order the outgoing edges of every vertex by
increasing key, breaking ties arbitrarily. Thus a lower lowpoint comes first,
and at an equal lowpoint a thin edge comes before a thick one. An edge with
empty $F(e)$ has lowpoint $h(v)$, and the constraint solver ignores it.

\paragraph{Computing the lowpoints.}
We compute these values without listing the sets $F(e)$. We keep a pair
$(\low,\lowtwo)$ for every edge. For a tree edge with source depth $c$, we
initialize both values to $c$. For a back edge $(v,u)$, we initialize them to
$(h(u),h(v))$. When an outgoing edge $f$ of $v$ is completed, we insert
$\low(f)$ and $\lowtwo(f)$ into the pair of the incoming tree edge of $v$, by
the following operation.
\begin{algorithmic}[1]
\Function{Insert}{$k$}
\If{$k<\low$}
\State $\lowtwo\gets\low$
\State $\low\gets k$
\ElsIf{$\low<k$ and $k<\lowtwo$}
\State $\lowtwo\gets k$
\EndIf
\EndFunction
\end{algorithmic}
This keeps the two smallest distinct values, with the source depth as a
repeated upper sentinel. Values at or above the sentinel are ignored. To prove
the recurrence for a tree edge $(u,v)$, split the back edges with tail in
$\subtree(v)$ into the back edges that leave $v$ and, for each child $w$ of
$v$, the back edges with tail in $\subtree(w)$. These parts belong to the
outgoing edges of $v$. The two smallest distinct depths below the sentinel in
a union are among the two smallest distinct depths below it in the parts: a
third distinct value of one part already has two smaller values in that part.
This proves the recurrence by induction in the postorder of the search. The correct thick bit
is
\begin{equation*}
\thick(e)=\bigl[\lowtwo(e)<h(\text{source of }e)\bigr].
\end{equation*}
Returns to the source of $e$ itself cannot set this bit, as required, because
they are not strict returns.

\paragraph{Sorting by key.}
We bucket-sort all edge records by their integer keys in one global pass, and
then append them to the outgoing lists of their sources. Keys lie in $[0,2N]$,
so one global scan of that range is enough. We do not allocate and scan a
bucket array of length $N$ separately for each vertex.

\paragraph{The search with arrays.}
A depth-first search with arrays needs no edge dictionary. When it scans an
incidence $v$--$w$:
\begin{itemize}[nosep]
\item if $w$ is unvisited, $w$ becomes a tree child of $v$;
\item otherwise, if $w$ is the parent of $v$, we skip this reverse tree
incidence;
\item otherwise, we create a back-edge record exactly when $h(w)<h(v)$.
\end{itemize}
The other incidences are ignored. Since there are no cross edges, every
undirected edge gets exactly one oriented record. Since the graph is simple,
skipping by parent vertex skips only the tree edge. Explicit frames
$(\text{vertex},\text{next adjacency index})$ produce the postorder of the
search without recursion.

\section{The constraints the algorithm decides}

Give every back edge one of two colors, $\cL$ and $\cR$. For distinct outgoing
edges $e,f$ of a vertex, write
\begin{equation*}
I(e,f)=\{b\in F(e):\ d(b)>\low(f)\}.
\end{equation*}
The \emph{fork constraints} are:
\begin{itemize}[nosep]
\item the set $I(e,f)$ is monochromatic;
\item the set $I(f,e)$ is monochromatic;
\item if both sets are nonempty, their colors are different.
\end{itemize}
An empty or one-element set is monochromatic and adds no equality. These are
the constraints of the left-right criterion: $I(e,f)$ is the set
$\mathrm{Interlaced}(e,f)$ of \cite{dFOMR06} (Theorem~2 on page~8 of the arXiv
version).

For a tree edge $p=(u,v)$, define its \emph{weak returns}
\begin{equation*}
W(p)=\{\text{back edges }(x,t):\ x\in\subtree(v),\ h(t)\le h(u)\}.
\end{equation*}
Equivalently, $W(p)$ is the union of the sets $F(e)$ over the outgoing edges
$e$ of $v$. We also require that all members of minimum depth of a nonempty
set $W(p)$ have the same color. We call this \emph{weak alignment}. The source
$u$ is included in the definition, so alignment is imposed before the returns
to $u$ are removed.

Parts~\ref{part:ii} to~\ref{part:iv} prove that these constraints, including
weak alignment, can be satisfied exactly when the graph is planar.

\section{Representation}

A \emph{list} consists of nodes for back edges, with a first pointer and a last
pointer. A node contains its depth key $d$ and a next pointer. Lists are ordered
by nonincreasing $d$. Concatenation is destructive: we connect the old last
node to the new first node, update the last pointer, and empty the header of
the list that was appended. This costs $O(1)$, whatever the number of nodes.

A \emph{block} $(L,R)$ consists of two such lists, not both empty. It means
that each list is monochromatic and, if both lists are nonempty, their colors
differ. Different blocks can be flipped independently.

A nonempty \emph{summary} is a triple $(a,A,P)$, where
\begin{itemize}[nosep]
\item the number $a$ is the least live depth;
\item the \emph{anchor} $A$ is a nonempty list that contains \emph{all} nodes of
depth $a$;
\item the \emph{block stack} $P$ is a stack of blocks, and its blocks contain
every other live node exactly once.
\end{itemize}
The anchor is one independent monochromatic class. Every node in $P$ has depth
strictly greater than $a$. For two blocks of the stack, every depth in the
upper block is at least every depth in the lower block. Equality is allowed.

For a nonempty list, its maximum is the key of its first node and its minimum
is the key of its last node. Let $\maxkey(\varnothing)=-1$. The maximum and
minimum of a block use the maxima and minima of its nonempty lists, so both
are $O(1)$ operations. The stack has a top pointer, and each block has a
pointer to the block below it. The anchor is stored separately, so we never
need to delete at the bottom of a singly linked stack.

\section{Merge pseudocode}

The operation \pr{Merge} transfers ownership destructively: it consumes $T$
and returns the updated $S$.
\begin{algorithmic}[1]
\Function{Merge}{$S,T$}
\If{$T$ is empty} \Return $S$ \EndIf
\If{$S$ is empty} \Return $T$ \EndIf
\State $a\gets S.a$
\State $b\gets T.a$ \Comment{the edge order guarantees $a\le b$}
\State $R\gets$ empty list
\While{$T.P$ is not empty}
\State $Q\gets\textsc{pop}(T.P)$
\If{$Q.L$ and $Q.R$ are both nonempty} \textbf{reject} the graph \EndIf
\State append the nonempty side of $Q$ to $R$
\EndWhile
\If{$b=a$}
\State append $T.A$ to $S.A$
\Else
\State append $T.A$ to $R$
\EndIf
\State $L\gets$ empty list
\While{$S.P$ is not empty and $\maxkey(\textsc{top}(S.P))>b$}
\State $Q\gets\textsc{pop}(S.P)$
\If{$\maxkey(Q.R)>b$} swap $Q.L$ and $Q.R$ \EndIf
\If{$\maxkey(Q.R)>b$} \textbf{reject} the graph \EndIf
\State append $Q.L$ to $L$
\State append $Q.R$ to $R$
\EndWhile
\If{$L$ or $R$ is nonempty} push $(L,R)$ onto $S.P$ \EndIf
\State \Return $S$
\EndFunction
\end{algorithmic}

This operation chooses no path, no embedding and no tentative fixed coloring.
Both rejections come from actual parity contradictions. Part~\ref{part:ii}
proves that no satisfiable instance is rejected.

\section{Trim pseudocode}

The operation \pr{Trim} deletes the returns to depth $c$.
\begin{algorithmic}[1]
\Function{Trim}{$S,c$}
\If{$S$ is empty} \Return empty \EndIf
\If{$S.a=c$} \Return empty \EndIf
\State \Comment{otherwise $S.a<c$, and all live depths are at most $c$}
\While{$S.P$ is not empty and $\minkey(\textsc{top}(S.P))=c$}
\State $\textsc{pop}(S.P)$
\EndWhile
\If{$S.P$ is not empty}
\State $Q\gets\textsc{top}(S.P)$
\State remove all first nodes of $Q.L$ whose depth is $c$
\State remove all first nodes of $Q.R$ whose depth is $c$
\State update the last pointer of an emptied list to null
\EndIf
\State \Return $S$
\EndFunction
\end{algorithmic}

All nodes in a popped block have depth $c$. At most one block can be partly
trimmed, and that block keeps at least one node of smaller depth. The anchor is
unchanged unless the entire summary becomes empty.

\section{Main loop}

We process the vertices in the postorder of the search. At a vertex $v$, we
start with $S$ empty and process the outgoing edges of $v$ in the order of
their keys:
\begin{itemize}[nosep]
\item for a back edge $(v,u)$, let $T$ be the one-node summary
$\bigl(h(u),[(v,u)],\text{empty block stack}\bigr)$;
\item for a tree edge $(v,w)$, let $T$ be the summary of $w$ computed earlier;
\item set $S\gets\pr{Merge}(S,T)$.
\end{itemize}
After all outgoing edges of $v$, if $v$ has a parent $u$, we store
$\pr{Trim}(S,h(u))$ as the summary of $v$. A root has an empty final summary.
If no merge rejected, we return true. The empty graph and isolated vertices
are accepted by the same code.

%======================================================================
\part{Complete proof of the constraint solver}\label{part:ii}

This part proves that the algorithm of Part~\ref{part:i} accepts exactly when
the fork constraints and weak alignments of Part~\ref{part:i}, Section~2 can be
satisfied together. The idea is that a summary is a normal form of a Boolean
formula. It describes exactly which colorings of the back edges that are still
active extend to colorings of all back edges below that satisfy every
constraint seen so far.

\subsection*{Representation refinement}

In this part a logical stack includes a virtual bottom block $(A,\varnothing)$
for the anchor. The physical representation of Part~\ref{part:i} stores the
anchor separately: the logical stack of $(a,A,P)$ is the virtual anchor
followed by the physical blocks of $P$. Every reference below to the bottom
block means this logical stack. \pr{Merge} never pops the virtual anchor,
because $a\le b$. It concatenates anchors directly when the minima agree.
\pr{Trim} removes the virtual anchor only when every active depth equals the
trimmed depth. Thus one-sided logical blocks and physical summary objects agree
exactly, and virtual anchors add no physical pops.

\subsection*{Definitions}

A back edge $f$ is a variable with key $d(f)$, the depth of its head. For an
outgoing edge $e$ of $v$, the set $F(e)$ consists of the back edges in the
branch of $e$ whose head is a strict ancestor of $v$, as in Part~\ref{part:i},
Section~1. Recall that $\low(e)$ is the least key in $F(e)$, with $h(v)$ as
sentinel when $F(e)$ is empty, and that $e$ is chordal when $F(e)$ has at least
two distinct keys. The outgoing edges are ordered by $(\low,\chordal)$, so a
thin edge comes before a chordal edge at equal lowpoint. Edges with empty $F(e)$
can be skipped in the constraint phase.

For a tree edge $(u,v)$, the set $W(u,v)$ is the union of $F(e)$ over the
outgoing edges $e$ of $v$. These are the weak returns from $\subtree(v)$, with
head depth at most $h(u)$. Weak alignment says that all members of $W(u,v)$
with minimum key have the same color. It is imposed before the key $h(u)$ is
deleted.

A block $(L,R)$ represents the formula: $L$ is monochromatic, $R$ is
monochromatic, and the colors of $L$ and $R$ differ when both are nonempty.
Thus a one-sided block is a freely colored monochromatic class, and every
nonempty block has exactly two colorings. There are no constraints between
different blocks. The lists $L$ and $R$ are singly linked and ordered by
nonincreasing key.

\subsection*{The invariant}

At a vertex $v$, the \emph{active} variables after a prefix of the outgoing
branches has been processed are the members of $F(e)$ for the processed
branches $e$. The stack invariant after processing such a prefix is the
following.
\begin{enumerate}[label=(I\arabic*),nosep]
\item Its lists are disjoint and cover exactly the active variables.
\item Its block formula is \emph{exactly} the existential projection, onto the
active variables, of the conjunction of: all fork constraints at vertices
inside the processed branches, all completed weak alignments (those of the
tree edges in the processed branches), all fork constraints at $v$ between
processed outgoing edges, and the weak alignment of the current minimum-key
active variables.
\item If a block $P$ is below a block $Q$, every key of $P$ is at most every key
of $Q$. The inequality is not strict.
\item When the stack is nonempty, its unique bottom block is one-sided and
contains exactly all minimum-key active variables. No other block contains that
minimum. We call it the anchor.
\end{enumerate}
The minimum key of every nonempty prefix is the minimum $a$ over all nonempty
outgoing branches, because the branches are processed by increasing lowpoint.
Therefore the partial alignment in (I2) is part of the final weak alignment of
the incoming tree edge of $v$, not an extra constraint.

\subsection*{Why completed formulas compose}

Back-edge variables of different outgoing branches are disjoint. All
constraints at $v$ involve only strict returns of those branches, and the
alignment at $v$ involves only the same active variables. So the hidden
variables of a branch, which were projected away, never occur in new
constraints. Formally, let $X$ and $Y$ be disjoint sets of hidden variables,
let $\Phi_{\mathrm{old}}$ be a formula on $X$ and the old active variables, let
$\Phi_{\mathrm{child}}$ be a formula on $Y$ and the child's active variables,
and let $\Psi$ be a formula on active variables only. Then
\begin{equation*}
\bigl(\exists X,Y.\ \Phi_{\mathrm{old}}\wedge\Phi_{\mathrm{child}}\wedge\Psi\bigr)
\iff
\bigl(\exists X.\ \Phi_{\mathrm{old}}\bigr)\wedge
\bigl(\exists Y.\ \Phi_{\mathrm{child}}\bigr)\wedge\Psi,
\end{equation*}
because $X$ occurs only in the first conjunct and $Y$ only in the second. This
lets us work only with the two block normal forms.

\subsection*{Exact aggregation of new fork constraints}

Suppose that $e_i$ is the new branch, with $i\ge2$, that $a$ is the lowpoint of
the first branch $e_1$, and that $b=\low(e_i)$. Thus $a\le\low(e_j)\le b$ for
$j<i$. Define
\begin{align*}
B&=\{f\in F(e_i):\ d(f)>a\},\\
U&=\Bigl\{f\in\textstyle\bigcup_{j<i}F(e_j):\ d(f)>b\Bigr\}.
\end{align*}
The newly added fork constraints, together with the earlier constraints, are
equivalent to:
\begin{equation*}
B\text{ monochromatic},\qquad U\text{ monochromatic},\qquad B\text{ opposite }U.
\end{equation*}
Here an empty set has its usual vacuous meaning.

\begin{proof}
For each $j<i$, write $M_j=\{f\in F(e_j):d(f)>b\}=I(e_j,e_i)$ and
$N_j=\{f\in F(e_i):d(f)>\low(e_j)\}=I(e_i,e_j)$. The new fork constraints are:
$M_j$ is monochromatic, $N_j$ is monochromatic, and $M_j$ is opposite $N_j$.
Since the first lowpoint is $a$, we have $N_1=B$, so the fork constraints force
$B$ to be monochromatic.

If $M_j$ is nonempty, then $N_j$ is nonempty. Indeed, if $\low(e_j)<b$, every
return of $e_i$ belongs to $N_j$. If $\low(e_j)=b$, then a nonempty $M_j$ means
that $e_j$ is chordal, and the order forces $e_i$ to be chordal too, so $e_i$
has a return strictly above $b$. Also $N_j\subseteq B$, since
$\low(e_j)\ge a$. Consequently every nonempty $M_j$ is opposite $B$, which
proves that all of $U$ is monochromatic and opposite $B$.

Conversely, $M_j\subseteq U$ and $N_j\subseteq B$, so the three aggregated
conditions imply every new fork constraint. This proves the equivalence. It
also shows that a nonempty $U$ forces $B$ to be nonempty.
\end{proof}

\subsection*{Child half of merge}

If $b=a$, we detach the child's anchor, which contains exactly its key-$a$
variables, and append its single list to the old key-$a$ anchor. This imposes
exactly the new weak alignment equalities. The detached anchor is independent
of every other child block, by (I2) and (I4). If $b>a$, we detach nothing. In
either case, the remaining child variables are exactly $B$.

Every remaining child block must be one-sided. A two-sided block already
requires opposite colors, so it cannot satisfy ``$B$ monochromatic''. If a
two-sided block occurs, rejection is therefore exact. Otherwise, we concatenate
the nonempty lists of these blocks from top to bottom into one list $R$.
Because the key ranges of the blocks are ordered by (I3), the list $R$ is still
nonincreasing. It contains exactly $B$ and represents exactly the old child
constraints together with ``$B$ monochromatic''. Every key in $R$ is at least
$b$, and strictly greater than $b$ when $b=a$.

\subsection*{Old half of merge}

We pop old blocks while their maximum key is greater than $b$. These are
exactly the old blocks that meet $U$: once a block with maximum at most $b$ is
reached, (I3) implies that all lower blocks also have maximum at most $b$.

For each popped block:
\begin{itemize}[nosep]
\item If both sides have a member of $U$, we reject. The existing constraint
that the two sides have opposite colors contradicts ``$U$ monochromatic''.
\item Otherwise exactly one side meets $U$. We swap the sides if needed, so that
it is $L$. Its entire list must have the common color of $U$, because that side
was already monochromatic. Its opposite list $R_{\mathrm{old}}$ has the
opposite color, and hence the color of $B$, since a nonempty $U$ forces $B$ to
be nonempty.
\end{itemize}
We concatenate these $L$ lists from top to bottom into a new list $L$. We
append each nonempty opposite list after the new child list $R$. We push
$(L,R)$ if it is nonempty.

This encodes \emph{exactly} the old block constraints together with ``$U$
monochromatic'' and ``$B$ opposite $U$''. Every side of an affected block is
forced into the stated class, and an unaffected block has no variable of $U$
and stays independent. No arbitrary extra equality is added.

\subsection*{List ordering and stack ordering}

The new list $L$ is nonincreasing, because the old blocks were popped from
higher keys to lower ones. Every old opposite list has all keys at most $b$,
whereas every key of the child list $R$ is at least $b$, so appending keeps the
order. In fact at most the \emph{lowest} popped block can have a nonempty
opposite side. If an upper popped block had a key at most $b$, then (I3) would
force every key of a lower block to be at most $b$, contradicting the fact that
the lower block was popped. So opposite lists are never interleaved.

When $U$ is empty, every old key is at most $b$. So the new one-sided block
$B$, whose minimum is at least $b$, can be pushed above all old blocks. When
$U$ is nonempty, every key left in the old stack is at most $b$, by the
stopping condition, and at most every key of every popped block, by (I3). Hence
it is at most every key of $B$ and every old key put in the new block. This
proves (I3) after the push.

The old anchor has only the key $a\le b$, and it is never popped. All variables
outside it have key greater than $a$, by (I4), and the child's key-$a$ anchor
was appended to it when necessary. Therefore (I4) is preserved, including the
uniqueness of the minimum-key anchor.

All list ownership is transferred, not copied. So (I1) follows from the
disjointness of the variables of different branches and from the transfer of
whole blocks.

\subsection*{Trim at the parent}

After all branches of $v$ have been processed, every active return has head key
at most $h(u)$, where $u$ is the parent of $v$. We delete every variable with
key $h(u)$. Because of (I3), we first pop all entire blocks whose minimum is
that key. Then at most one remaining block can contain that key, and in each of
its lists these variables form a prefix. We delete the prefixes, and delete the
block if both sides become empty.

This is an \emph{exact} existential projection of a block formula. To extend a
coloring of the survivors, give each deleted variable the color of the
survivors on its side. If one side or both sides have no survivors, give the
missing side the color opposite to the other side, or an arbitrary color when
no survivor remains. Thus (I2) is preserved, and now the weak alignment of
$(u,v)$ is a completed alignment. The remaining variables are exactly
$F((u,v))$. Ordered lists and the stack order persist under deletion. If
anything remains, the old minimum $a$ is less than $h(u)$, so its untouched
anchor is still the unique minimum anchor. Otherwise the stack is empty.

\subsection*{Base and completion}

A direct outgoing back edge gives a one-sided block with one variable. It has
no internal constraints and satisfies all invariants. A processed tree edge to a
child supplies the invariant by induction. At a root of the search there are
no strict returns, so all variables have been existentially eliminated. A
successful empty normal form therefore means that all constraints in that
component have a satisfying assignment. Rejection happens exactly when an
explicitly required equality would identify two opposite nonempty classes. By
Part~\ref{part:iii} and Part~\ref{part:iv}, the algorithm accepts exactly when
the graph is planar.

\subsection*{Linear work}

We make one list node per back edge, and singly linked stacks of block records.
We store list heads and tails, block minima and maxima, and anchor pointers. A
one-edge summary creates one block. Each merge of an outgoing edge creates at
most one new block and permanently consumes every popped block of the child and
of the old stack. A discarded anchor block is consumed as well. Thus $O(m+n)$
blocks are created and $O(m+n)$ block operations occur, where $n$ and $m$ are
the numbers of vertices and edges. Each active back-edge node is deleted once,
so the total work of \pr{Trim} is $O(m)$. Constant-time pointer
concatenations never scan a list. With arrays, explicit search frames and the
integer-key counting sort, the bound $O(n+m)$ is deterministic on a word RAM.

%======================================================================
\part{Weak alignment does not change satisfiability}\label{part:iii}

The algorithm decides the fork constraints together with all weak alignments.
This part shows that adding the weak alignments does not change whether the
fork constraints can be satisfied.

Let $p=(u,v)$ be a tree edge, let $W(p)$ be its set of weak returns, and let
$M(p)$ be the set of its members of minimum depth $\ell$. Assume that $W(p)$ is
nonempty. We read the fork constraints as pairwise equalities within each
nonempty set $I(e,f)$ and pairwise inequalities between the two nonempty sets.
Only distinct variables give a nontrivial equality. We call a variable
\emph{isolated} if it occurs in no nontrivial equality or inequality, and
\emph{nonisolated} otherwise.

Suppose that $b\in M(p)$ occurs in one of these equalities or inequalities, at
a fork $w$. For $b$ to belong to $I(e,f)$, the other branch $f$ must have
lowpoint below $\ell$. Consequently $f$ has a back edge with head depth below
$\ell$. The vertex $w$ lies on the tree path from the root to the tail of $b$.
It cannot belong to $\subtree(v)$. If it did, the back edge of $f$ would have
its tail in $\subtree(v)$ and its head depth below $\ell\le h(u)$,
contradicting the choice of $\ell$ as the minimum depth in $W(p)$.

Therefore $w$ is a strict ancestor of $v$. All members of $M(p)$ then belong to
the same outgoing branch $e$ of $w$, the one containing $v$, and have the same
head depth $\ell$. So they all belong to the same set $I(e,f)$, and the
original constraints already force them to have one color.

It follows that each $M(p)$ either is already forced to be monochromatic by
the original constraints, or consists entirely of variables that occur in no
nontrivial original equality or inequality. To satisfy all alignments at once,
we keep a satisfying coloring of every nonisolated variable and give every
isolated variable the color $\cR$. A set of the first kind is already
monochromatic, and a set of the second kind now is. The fork constraints still
hold, since no nonisolated variable changed its color. Thus the original fork constraints can be satisfied if
and only if the fork constraints together with \emph{all} weak alignments can
be satisfied.

%======================================================================
\part{The planarity characterization, proved constructively}\label{part:iv}

This part proves that the constraints of Part~\ref{part:i}, Section~2 can be
satisfied exactly when the graph is planar. We first show that a plane
embedding gives a coloring that satisfies the fork constraints. Then
Sections~1 to~5 show that a coloring that satisfies the fork constraints and
the weak alignments gives an explicit plane drawing. Sections~6 and~7 supply
the step from a plane embedding to a nonalternating contour, which the first
argument assumes.

\paragraph{Contours and nonalternation.}
A \emph{rotation} at a vertex is a cyclic order of its incident edges. Given a
rotation at every vertex, we build a word of tokens, the \emph{contour}. At a
nonroot vertex, we cut its rotation immediately after the parent incidence, so
that it becomes a linear order. At the root, we cut at any corner. Then we
replace each outgoing tree incidence, recursively, by the contour of the child.
A back edge $(x,u)$ has a \emph{source token} at $x$ and a \emph{target token}
at $u$. Two back edges \emph{alternate} if exactly one token of the second lies
strictly between the two tokens of the first. A contour is
\emph{nonalternating} if no two back edges alternate in it.

\paragraph{The fork constraints are necessary.}
Assume that the rotations of a plane embedding give a nonalternating contour.
Sections~6 and~7 prove this assumption. We declare a back edge $(x,u)$ to be
$\cL$ if its target token occurs before the child incidence at $u$ leading to
$x$, and $\cR$ otherwise.

Take two back edges $b_1=(x_1,u_1)$ and $b_2=(x_2,u_2)$ whose fundamental tree
paths share an edge, with $h(u_1)>h(u_2)$. The \emph{fundamental tree path} of
a back edge $(x,u)$ is the tree path from $u$ to $x$. Let $w$ be the lowest
common ancestor of $x_1$ and $x_2$, and let $e_1,e_2$ be the distinct outgoing
edges of $w$ that lead to the respective source tokens. Such an edge can itself
be an outgoing back edge, when its source is $w$. Both sources lie in the same
child subtree of $u_1$. The target token of $b_1$ occurs before both source
tokens when $b_1$ is $\cL$, and after both when it is $\cR$. The target token
of $b_2$ is outside the entire interval of the subtree of $u_1$. Checking the
order of these four tokens, nonalternation implies:
\begin{equation}\label{eq:forkorder}
e_1\text{ precedes }e_2\text{ at }w\quad\text{if and only if}\quad b_1\text{ is }\cL.
\end{equation}

Now take distinct outgoing edges $e,f$ of a vertex $w$. For any $b\in I(e,f)$,
choose a return $c\in F(f)$ of minimum depth. It exists because
$d(b)>\low(f)$, whereas $d(b)<h(w)$. The two fundamental paths share the
incoming edge of $w$, and the head of $b$ is deeper than the head of $c$. So
\eqref{eq:forkorder} gives the same color to every member of $I(e,f)$,
determined only by whether $e$ precedes $f$. Reversing the roles gives the
other color to every member of $I(f,e)$. These are exactly the fork
constraints.

Thus a plane embedding gives a coloring that satisfies the fork constraints,
and Part~\ref{part:iii} makes it weakly aligned. The following construction
proves the converse in full.

\section{Definitions}

Root the search tree $T$. All non-tree edges $b=(x,u)$ go from a descendant $x$
to a strict ancestor $u$, and $h(v)$ is the depth. For an outgoing edge $e$ of
$v$, the set $F(e)$ of its returns and the values $\low(e)$, $\chordal(e)$ and
$\key(e)=2\low(e)+\chordal(e)$ are as in Part~\ref{part:i}, Section~1. When
$F(e)$ is nonempty, let $\high(e)$ be the maximum head depth in $F(e)$.

We assume a coloring with the following properties. For each fork at a nonroot
vertex $v$, and distinct outgoing edges $e,f$ of $v$:
\begin{itemize}[nosep]
\item the set $\{b\in F(e):d(b)>\low(f)\}$ is monochromatic;
\item the analogous set with $e$ and $f$ exchanged is monochromatic;
\item if both sets are nonempty, their colors are opposite.
\end{itemize}
Moreover the coloring is weakly aligned: for every tree edge $(p,v)$, among the
back edges with source in $\subtree(v)$ and head equal to $p$ or an ancestor of
$p$, all those of minimum head depth have the same color. If this set is empty
there is no condition. The inclusion of $p$ is deliberate. It covers the
alignment just before the stack algorithm removes the returns to $p$.

We extend the colors to outgoing edges. For an outgoing tree edge $e$ with
$F(e)$ nonempty, let $\sigma(e)$ be the color of any return of depth $\high(e)$.
If $F(e)$ is empty, choose any color. For a back edge, $\sigma(e)$ is its given
color.

We order the outgoing edges at $v$ by putting all $\cL$ edges first, in
decreasing key, and then all $\cR$ edges, in increasing key. Ties are broken
arbitrarily. Every comparison used below is strict, so the tie breaking does
not matter.

We run a depth-first enumeration of these ordered outgoing edges, and give
every back edge a distinct \emph{source rank} when its outgoing occurrence is
reached. Incoming back-edge occurrences do not enter this enumeration.

\section{An explicit contour word and rotation}

At a vertex $v$, each outgoing tree edge $e=(v,w)$ receives two lists:
\begin{align*}
L_e&=\text{all back edges $(x,v)$ with $x\in\subtree(w)$, colored $\cL$},\\
R_e&=\text{the corresponding $\cR$ edges}.
\end{align*}
In each list, we order the edges by decreasing source rank.

We replace each outgoing tree edge $e=(v,w)$ by the block
\begin{equation*}
\text{target ports }L_e;\quad\text{child }w;\quad\text{target ports }R_e,
\end{equation*}
and each outgoing back edge $b$ by its source port. At nonroot vertices, the
incoming tree edge precedes this list in the cyclic rotation. These rules give
a rotation at every original vertex.

Equivalently, add a distinct leaf of degree one for every source port and every
target port. The resulting rooted ordered tree $T^*$ has a leaf contour word
$W$. Inserting target ports does not change the order of the source ports, so
the source ranks used to build the lists are exactly the order of the source
ports in $W$.

For a tree edge $e=(v,w)$, its whole block $L_e\,W(w)\,R_e$ is contiguous,
where $W(w)$ is the contour of the subtree of $w$. More generally, $W(w)$ is a
contiguous subword of $W$ for every rooted subtree. At the lowest common
ancestor $v$ of two sources, the order of their source ports is the order of
their distinct outgoing fork edges at $v$. If a source is $v$ itself, we use
the corresponding back edge as that outgoing fork edge.

\section{The fork-order lemma}

Consider back edges $b_1=(x_1,u_1)$ and $b_2=(x_2,u_2)$ whose tree paths share
an edge. Let $v$ be the lowest common ancestor of $x_1$ and $x_2$, and let
$e_1,e_2$ be their distinct outgoing fork edges at $v$. Both $u_1$ and $u_2$
are strict ancestors of $v$.

\begin{lem}{(A)}
Suppose that $h(u_1)>h(u_2)$. Then $e_1$ precedes $e_2$ in the outgoing order
if and only if $b_1$ is $\cL$.
\end{lem}

\begin{proof}
Since $h(u_1)>h(u_2)\ge\low(e_2)$, the first fork constraint makes all returns
of $e_1$ above $\low(e_2)$ have the color of $b_1$. In particular its highest
return does, so $\sigma(e_1)$ is the color of $b_1$. If $\sigma(e_2)$ has the
other color, the conclusion follows from the rule that $\cL$ edges come before
$\cR$ edges. Otherwise both edge colors are the same. Then the two opposite
fork sets cannot both contain their highest returns, so
\begin{equation*}
\high(e_2)\le\low(e_1).
\end{equation*}
Consequently $\low(e_2)\le\high(e_2)\le\low(e_1)$. If
$\low(e_2)<\low(e_1)$, then $\key(e_2)<\key(e_1)$. If the lowpoints are equal,
every return of $e_2$ is at that common lowpoint, whereas $b_1$ returns
strictly higher. Thus $e_2$ is not chordal and $e_1$ is chordal, which again
gives $\key(e_2)<\key(e_1)$. The signed order therefore puts $e_1$ first
exactly when $b_1$ is $\cL$.
\end{proof}

\begin{lem}{(B)}
Suppose that $u_1=u_2=u$, that $b_1$ is $\cL$ and that $b_2$ is $\cR$. Then
$e_1$ precedes $e_2$.
\end{lem}

\begin{proof}
Put $l_i=\low(e_i)$, so $l_i\le h(u)$.
\begin{itemize}[nosep]
\item If both $l_i<h(u)$, the fork constraints force $\sigma(e_1)=\cL$ and
$\sigma(e_2)=\cR$.
\item If $l_1<h(u)=l_2$, they force $\sigma(e_2)=\cR$, and
$\key(e_1)<\key(e_2)$. Whether $\sigma(e_1)$ is $\cL$ or $\cR$, the edge $e_1$
precedes $e_2$.
\item If $l_2<h(u)=l_1$, then symmetrically $\sigma(e_1)=\cL$ and
$\key(e_2)<\key(e_1)$, so again $e_1$ precedes $e_2$.
\item Suppose that $l_1=l_2=h(u)$. Let $(p,v)$ be the incoming tree edge of $v$.
If its weak minimum head depth were $h(u)$, alignment would force $b_1$ and
$b_2$ to have the same color. Therefore that weak minimum is strictly below
$h(u)$. Since $h(p)\ge h(u)$, both $b_1$ and $b_2$ are weak returns of $(p,v)$,
so some third outgoing edge $e_0$ of $v$ has $\low(e_0)<h(u)$. The fork
constraints with $e_0$ then make every return of $e_i$ have the color of $b_i$,
including its highest return. Thus $\sigma(e_1)=\cL$ and $\sigma(e_2)=\cR$,
which proves the claim.\qedhere
\end{itemize}
\end{proof}

\section{No two back-edge port pairs alternate in \texorpdfstring{$W$}{W}}

Let $b_i=(x_i,u_i)$, for $i=1,2$, be distinct back edges.

\paragraph{Paths that share no edge.}
If their fundamental tree paths share no tree edge, their port pairs do not
alternate. Here are the complete structural cases. If the heads are
incomparable in the rooted tree, both pairs occur in disjoint subtree words.
Otherwise we may assume that $u_1$ is an ancestor of $u_2$. If $x_1$ lies
outside $\subtree(u_2)$, both ports of $b_2$ lie in $W(u_2)$, which contains
neither port of $b_1$. If $x_1$ lies in that subtree, but the two paths share
no edge, then either their sources diverge at $u_2$ or $x_1=u_2$. The ports of
$b_2$ both lie in the single outgoing child block of $u_2$ leading to $x_2$.
The source port of $b_1$ is outside that block, and so is its target port. If
$u_1=u_2$ and the sources diverge immediately, the pairs lie in disjoint
outgoing child blocks. In all cases one pair is in a contiguous subword that
contains neither port of the other, which excludes alternation.

\paragraph{Paths that share an edge.}
Now suppose that their paths share an edge. The heads are comparable, and both
are strict ancestors of the source fork $v$.

First suppose that $h(u_1)>h(u_2)$. The target port of $b_2$ lies outside
$W(u_1)$. The target port of $b_1$ lies on the $\cL$ side or the $\cR$ side of
the child block of $u_1$ that contains $v$. Both source ports lie inside that
block. If $b_1$ is $\cL$, its target comes before both source ports, and
alternation would require the source of $b_2$ to come before the source of
$b_1$. If $b_1$ is $\cR$, its target comes after both source ports, and
alternation would require the source of $b_1$ to come before the source of
$b_2$. Both possibilities contradict Lemma~(A), since the order of the source
ports is the order of $e_1$ and $e_2$ at $v$.

Finally suppose that $u_1=u_2$. If their colors agree, their target ports are
in one list on the same side, ordered in reverse order of their source ports.
Their four ports are therefore nested, not alternating. If their colors differ,
say $b_1$ is $\cL$ and $b_2$ is $\cR$, the target of $b_1$ precedes the shared
child block and the target of $b_2$ follows it. Alternation would require the
source of $b_2$ to come before the source of $b_1$, contrary to Lemma~(B). This
exhausts all cases.

\section{Convert a nonalternating contour to explicit curves}

This step applies to any finite rooted ordered tree with a matching on some of
its leaves whose pairs do not alternate in the leaf order. We apply it to the
tree $T^*$ of Section~2, where the matching pairs the two ports of each back
edge.

Add one dummy leaf to each original vertex which otherwise has no child or port
leaf. All leaves are now genuine ports or dummy ports. Let $D$ be their maximum
depth. Subdivide each leaf edge as necessary so that every leaf has depth $D$.
Subdivision does not change the leaf order. Number all leaves $1,\ldots,r$ in
order.

For an internal vertex $z$, let $L(z)$ and $R(z)$ be the smallest and largest
indices of its descendant leaves, and put
\begin{equation*}
\pos(z)=\Bigl(\frac{L(z)+R(z)}{2},\ D-\operatorname{depth}(z)\Bigr).
\end{equation*}
Put leaf $i$ at $(i,0)$, and draw every tree edge as a straight segment.

\paragraph{This is a plane tree drawing.}
The descendant-leaf intervals of distinct vertices at the same depth are
disjoint. Hence edges from distinct parents in one level strip lie in disjoint
horizontal intervals. Edges from the same parent end at distinct
$x$-coordinates on the next horizontal level, and their straight segments meet
only at the parent. Distinct open level strips are disjoint. Their common
boundary lines meet edges only at endpoints, and distinct vertices on one such
line have distinct $x$-coordinates. Thus tree-edge images meet only at common
endpoints. Every original vertex has a descendant leaf at strictly larger
depth, so all original vertices lie strictly above the line $y=0$.

\paragraph{The back edges.}
For each matched pair of genuine port leaves $i<j$, draw the polygonal U-shaped
path with three segments
\begin{equation*}
(i,0),\quad(i,-(j-i)),\quad(j,-(j-i)),\quad(j,0).
\end{equation*}
Distinct matching pairs have four distinct endpoints. Nonalternation says that
their intervals are either disjoint or one is strictly inside the other.
Disjoint intervals give disjoint horizontal ranges for the two U paths. If
$[k,l]$ is strictly inside $[i,j]$, the smaller U lies strictly between the
vertical sides of the larger U and strictly above its bottom, since $l-k<j-i$.
Thus the U paths are disjoint. All their interior points have $y<0$, so they
also miss the tree drawing, which lies in $y\ge0$. Every coordinate is
rational.

Each original non-tree edge is now represented by its source-port path, its
lower U path, and its target-port path. These three pieces form a simple curve
between the original endpoints. Port paths for distinct ports are disjoint
except when they are attached at the same original vertex, and the U pieces are
pairwise disjoint. Consequently distinct original edges meet exactly when they
have a common endpoint, and then only there. Delete the dummy port paths and
suppress all port vertices and subdivision vertices. This is an explicit plane
drawing of the original graph.

For a disconnected graph, construct a drawing of each component independently
and place translated and scaled copies in disjoint rectangles.

\section{Optional necessity by face permutations}

Suppose that a plane embedding is given combinatorially, as a rotation system
whose number of faces satisfies Euler's formula with characteristic $2$. Then
the step from the embedding to a nonalternating contour needs no topology. We
treat a connected graph. Contract the tree edges one at a time. In a dart
rotation, contracting an edge between distinct vertices deletes its two darts
and splices the two vertex cycles. For the face permutation
$\varphi=\text{rotation}\circ\text{reversal}$, this only deletes those darts
from their face cycles, so the number of faces is preserved. The single
exceptional component with no darts counts as having one face.

After all tree edges are contracted, the remaining rotation is exactly the
contour word $W$, with $b$ matched pairs, one vertex, and $b+1$ faces by
Euler's formula. If $b=0$, the word $W$ is empty and nonalternation holds
trivially. Assume that $b>0$. If $W$ has no cyclically adjacent matched pair,
then $\varphi$ has no cycle of length one, so it has at most $2b/2=b$ cycles,
a contradiction. Delete a cyclically adjacent matched pair. One of its darts
was a face of length one, and deleting both darts decreases the number of faces
by one. For $b=1$, the two faces of length one become the one empty face. By
induction, all pairs can be deleted in this way. An adjacent pair cannot
alternate with any other pair, and deletion preserves alternation among the
remaining pairs. Therefore $W$ is nonalternating.

This argument uses the exact Euler identity for rotation systems, not just the
inequality $m\le3n-6$. For embeddings given geometrically, Section~7 proves the
same step from the edge bound alone.

\section{Geometric necessity using only the Euler edge bound}

If the available fact from Euler is only $m\le3n-6$, the following argument
gives the step from a plane drawing to a nonalternating contour. It uses
elementary compactness and continuity, and finite polygonal paths. The
topology of a tree neighborhood is spelled out in Section~7b. We treat one
component of the graph and its search tree. For a disconnected graph, we apply
the argument to each component.

\subsection{Polygonalize a finite continuous plane drawing}

Parameterize every edge image by an injective continuous arc
$\gamma\colon[0,1]\to\mathbb R^2$. At each vertex $u$, choose a small closed
disk $D_u$ centered at $u$. We choose these disks pairwise disjoint, and we
require $D_u$ to avoid every edge image not incident to $u$. For an incident
edge parameterized from $u$, we also require $D_u$ to avoid
$\gamma([1/3,1])$. These requirements can be met, because they concern finitely
many compact sets that do not contain $u$. For $\gamma$ from $u$ to $v$, let
$t_u$ be the last parameter at which $\gamma$ lies on the boundary of $D_u$, and
$t_v$ the first parameter at which it lies on the boundary of $D_v$. Since the
same edge parameterized from $v$ is $t\mapsto\gamma(1-t)$, the disk $D_v$ avoids
$\gamma([0,2/3])$. Continuity gives $t_u<1/3<2/3<t_v$, and the middle arc
$\gamma([t_u,t_v])$ lies strictly outside the two disks except at its ends. It
avoids all other vertex disks as well.

Replace each end part by the radial segment from its vertex to its exit port.
Radial segments of different incident edges are disjoint except at their
vertex, because their exit ports are distinct points of one circle. Every
middle arc is disjoint from all radial segments of other edges. The middle arcs
are pairwise disjoint compact sets. Therefore we can choose for each middle arc
a small open metric neighborhood, so that these neighborhoods are pairwise
disjoint and avoid all radial segments of other edges and every original
vertex. Uniform continuity of the arc gives a finite partition whose
consecutive sample points can be joined by straight segments inside this
neighborhood: choose the samples so close that each such segment stays inside a
small ball about one sample point.

Append the two radial segments of the edge. This gives a finite polygonal walk
between the original ends, which misses every other constructed edge walk
except at shared original ends. To remove self-intersections, split its
segments at all isolated pairwise intersections and at the endpoints of
collinear overlaps, identify coincident subsegments, and take a simple path in
the resulting finite arrangement of segments. This is a simple polygonal arc
between the same ends. Doing this for each edge gives a polygonal plane
drawing. The argument uses neither the Jordan curve theorem nor the
Schoenflies theorem.

\subsection{A thin neighborhood of a polygonal embedded tree is a disk}

Subdivide the drawn tree at each bend, without changing its image. Around each
tree vertex, choose a small convex polygon. It has one facet perpendicular to
each incident edge direction, and the edge crosses the relative interior of
that facet. Such a polygon can be made by intersecting finitely many
half-planes tangent to a small circle, including one outward normal in each
required direction. Extra normals can be added to make the intersection
bounded. Distinct incident edges have distinct directions near their common
endpoint. We make the polygons disjoint, and small enough to avoid all edges
not incident to their vertex and all later segments of incident edges. Along
each straight tree segment between its end polygons, we place a thin closed
rectangular strip. It meets each end polygon exactly in a short boundary
interval of the right facet. Otherwise all strips and polygons have disjoint
interiors and no unintended meetings. Finitely many positive separation
distances allow all widths to be chosen. The strip intervals on one polygon
are pairwise disjoint.

The union $N$ is a closed disk. Here is an elementary gluing proof. Every
convex polygon is a disk, by radial parametrization. The union of two disks
that meet in exactly one nondegenerate proper closed boundary interval is a
disk. To see this, use the given disk parametrizations to map the two disks to
the two halves of a square, with the gluing interval going to the common side.
Choose boundary homeomorphisms that agree pointwise on the common interval, and
extend them over the complementary boundary arcs. Every homeomorphism $f$ of
the circle extends explicitly over the unit disk by $H(0)=0$ and $H(rz)=rf(z)$
for $|z|=1$ and $0<r\le1$. Its inverse is the radial extension of $f^{-1}$.
Radial parametrizations of the convex half-squares give the same construction
for those domains. Thus the two disk maps can be chosen to agree on the common
side. Glue the half-squares to one square. The induced map from that square to
the union is continuous and bijective, hence a homeomorphism, because its
domain is compact and the plane is Hausdorff. The outer boundary of the square
maps to the boundary formed by the two old boundaries without the interior of
the glued interval. This boundary property is kept at each attachment. For the
unions of polygons and strips here, this is also their boundary in the plane,
as one checks locally at points of a facet, of a corner, or of a glued
interval. Starting from one vertex polygon, we attach one tree strip and then
its new vertex polygon at a time. Each attachment meets the previous union in
exactly one boundary interval, so the gluing lemma applies by induction.

All images of non-tree edges, except their end tails, lie outside $N$, and each
meets the boundary exactly once at each end. Traversing the boundary of $N$
visits these ports in exactly the contour order given by the cyclic orders of
the edges at each original vertex. This follows by induction from the same
attachment of leaves: a strip substitutes the boundary contour of the child for
the interval where it was attached, which is exactly the recursion of the
rooted-tree contour.

\subsection{Alternating exterior arcs would give a plane \texorpdfstring{$K_5$}{K5}}\label{sec:k5}

Suppose that two non-tree edges have alternating boundary ports
$p_1,p_2,p_3,p_4$, with exterior arcs joining $p_1$ to $p_3$ and $p_2$ to
$p_4$. Use the four boundary arcs between consecutive ports as a $4$-cycle. A
homeomorphism from the unit disk onto $N$ gives one interior point $z$ and four
internally disjoint arcs from $z$ to the four ports: the images of the four
radial segments. These four spokes are inside $N$, and the two exterior arcs are
outside $N$. Consequently the cycle, the four spokes and the two exterior arcs
form a plane drawing of $K_5$. It has five vertices and ten edges,
contradicting Euler's bound $10\le3\cdot5-6=9$. Therefore no two pairs of
ports of non-tree edges alternate.

This proves that a plane drawing gives a nonalternating tree contour, using
only the Euler edge bound as a fact about planar graphs. Section~5 proves the
reverse implication by explicit coordinates. So the proof works both with
geometric embeddings and with rotation systems, and no fast planarity
algorithm is assumed in either direction.

%======================================================================
\part{Word RAM bound and use in the crossing extractor}\label{part:v}

This part shows that the algorithm is one fixed word-RAM program with a linear
bound, collects the correctness proof, and explains how the crossing extractor
uses it.

\section{A single fixed program}

We use vertex labels $0,\ldots,N-1$, edge identifiers $0,\ldots,M-1$, and null
pointers. All working structures are arrays of records of fixed size: search
frames, depth and parent arrays, two lowpoint arrays, bucket links, links of
outgoing lists, back-edge nodes, summary records, and block records. Fields
contain only indices, depths in $[0,N]$, or Boolean values. Array capacities
are allocated from $N$ and $M$. If $M$ is not part of the input header, we
first count the adjacency entries in time $O(N+M)$. The word size need only
address $O(N+M+1)$ words.

Every line of the pseudocode expands into a fixed number of RAM operations:
read or write a word, compare, branch, add or subtract an index, or move a
pointer. The multiplicative factor does not depend on the input graph. Appending
to a linked list changes at most a constant number of fields. No implicit
copying of a list or a stack is permitted or needed.

\section{Total operations}

The first search scans $2M$ adjacency entries and has $N$ vertex-completion
steps. Every edge causes at most two lowpoint insertions into the record of
its parent edge. The global bucket pass scans $O(N)$ buckets and $M$ edge
records. The postorder constraint pass visits $N$ vertices and $M$
outgoing-edge records.

One return-list node is allocated per back edge, so at most $M$ nodes in total.
Only one-edge summaries of back edges allocate summary objects. All returns of
child tree edges transfer the ownership of an existing object. Each effective
merge creates at most one new block, so there are at most $M$ block creations.
Every popped block is permanently consumed: it is never pushed again. Thus all
child pop loops, all pop loops of the previous stack and all whole-block trim
loops together execute at most $M$ block pops. Every iteration of a partial
trim removes one return node permanently, so these iterations total at most
$M$. Dropping a whole anchor or a whole block does not traverse its nodes. Each
tree edge has one call of \pr{Trim}, with at most one partly trimmed block.

Consequently the total number of RAM instructions is bounded by
\begin{equation*}
C\,(N+M+1)
\end{equation*}
for a fixed constant $C$, and the total number of allocated words has the same
form. This is a deterministic worst-case bound. It assumes neither average-case
hashing nor amortized union-find operations. The charging argument for linked
lists above bounds the sum of the operations in one run, in the worst case.

\section{Correctness theorem}

The two implications now have explicit finite chains of reasons.

If the graph has a plane embedding, Part~\ref{part:iv}, Section~7 (or the
rotation proof in Section~6) gives a nonalternating contour for the search
tree. The fork comparison at the beginning of Part~\ref{part:iv} gives a
coloring that satisfies every fork constraint. Part~\ref{part:iii} gives a
weakly aligned one. The exact projection invariant of Part~\ref{part:ii}
implies that no merge rejects. Hence the program returns true.

If the program returns true, the final empty summary of every component is a
satisfiable existential projection. Reversing the finite projections and
merges, using their extension properties proved in Part~\ref{part:ii}, gives a
full aligned coloring. Part~\ref{part:iv}, Sections~2 to~4 produce a
nonalternating contour. Section~5 gives an explicit rational polygonal drawing
of that component. We put distinct components in disjoint translated and scaled
rectangles. Thus the graph has a plane embedding.

The coloring and the drawing are witnesses used only in the correctness proof.
The Boolean program does not construct them, so their construction costs
nothing in the linear bound.

\section{Use in the crossing extractor}

The crossing extractor of the companion note \cite{HJ26} calls a Boolean
planarity test inside its classification step, and states the test as Fact~P.
The program of this note is such a test, and the extractor uses it without any
change to its classification step. Each planarity call still costs
$O(|V(K)|+|E(K)|+1)$ on the current graph $K$. All bounds of the companion
note on the number and the sizes of the classification calls therefore remain
valid, and the crossing extractor still runs in time $O(nm)$.

The planarity program uses one full depth-first search and linear scans. It
calls no path search and no flow computation between two vertex sets. Thus it
adds no searches to the search count of the companion note, beyond constant
work per planarity call.

\section{Tests}

The file \texttt{planarity\_linear.py} in this folder implements the program of
Part~\ref{part:i} in Python, with arrays and linked records as in
Part~\ref{part:i}, Section~3. Its function \texttt{is\_planar} takes the
adjacency lists of a simple graph. We compared it with the independent
left-right implementation \texttt{check\_planarity} of NetworkX~3.2.1
\cite{HSS08}. The answers agreed on all $33{,}868$ simple labeled graphs with
at most six vertices, including the empty graph, and on $10{,}000$ random
graphs with up to $150$ vertices and random orders of the adjacency lists, of
which $5{,}910$ are planar. These are checks of the implementation. The proof
above, not these finite tests, establishes correctness for every input. Python
lists are not the arrays of a word RAM, so the implementation illustrates the
program but does not prove its time bound.

%======================================================================
\part{The Lean formalization}\label{part:vi}

The Lean formalization of this note is in the repository \texttt{GraphMinor}.
A plane embedding is the structure \texttt{HP25.PlaneEmbedding} of the
formalization of Humeau and Pous: the vertices are distinct points of the
complex plane, and the edges are arcs that meet only at common ends.
\begin{itemize}[nosep]
\item The folder \path{GraphMinor/LRPlanarity} proves the mathematics of
Parts~\ref{part:i} to~\ref{part:iv}. Its file \path{Statements.lean} lists the
statements, each with the place of the note it formalizes. The search and the
lowpoints follow Part~\ref{part:i}, Section~1. The solver follows
Part~\ref{part:ii} (\path{Merge.lean}, \path{Trim.lean},
\path{Aggregation.lean}, \path{Solver.lean}), weak alignment follows
Part~\ref{part:iii} (\path{Alignment.lean}), and the files
\path{ForkOrder.lean}, \path{Contour*.lean}, \path{TreeDrawing*.lean},
\path{Polygonalization*.lean}, \path{TreeNeighborhood*.lean} and
\path{ExteriorK5*.lean} follow Part~\ref{part:iv}.
\item In the folder \path{GraphMinor/Compute}, the files
\path{PlanarityRam*.lean} give the word-RAM program and prove its step bound,
following Part~\ref{part:v}, Sections~1 and~2. The file
\path{SpecsPlanarity.lean} states what the program computes.
\item The statement \texttt{PlanarityTestStatement}, in the file
\path{CitedPlanarity.lean} of the same folder, is Fact~P of the companion
note. It is proved with no hypotheses, as
\texttt{planarityTestStatement\_proved} in \path{DisjointPathsFinal.lean}, by
assembling the parts as in Part~\ref{part:v}, Section~3. The same file then
proves the main result of Graph Minors XIII,
\texttt{disjointPathsStatement\_proved}.
\end{itemize}

The Lean formalization differs from this note in the following places.
\begin{itemize}[nosep]
\item \emph{The model of computation.} The programs are fixed word-RAM
commands with a logarithmic word size, and every read, write and allocation
has a proved bound. A null pointer is zero, an optional index is stored plus
one, and pointer arenas reserve slot zero. The step count uses adjacency
entries rather than undirected edges, which changes only the constant.
\item \emph{The coordinates of Part~\ref{part:iv}, Section~5.} Leaves are
numbered from zero. One leaf word numbers the leaves of all components, and each
component lies in its own vertical strip, instead of in a separate translated
and scaled rectangle. The segment to a port is followed by a vertical segment
that contains the subdivision vertices and the vertical side of its U path.
After the subdivision vertices are suppressed, these are the same drawings.
\item \emph{The route to necessity.} The Lean works with geometric embeddings
and follows Part~\ref{part:iv}, Section~7. It does not use the argument of
Section~6 for rotation systems.
\end{itemize}

{\small
\begin{thebibliography}{dFOMR06}

\bibitem[Br09]{Br09}
U.~Brandes,
The left-right planarity test,
manuscript, Universit\"at Konstanz, 2009.

\bibitem[dFOMR06]{dFOMR06}
H.~de~Fraysseix, P.~Ossona de Mendez and P.~Rosenstiehl,
Tr\'emaux trees and planarity,
\emph{Internat. J. Found. Comput. Sci.} \textbf{17} (2006), 1017--1029.
Preprint arXiv:math/0610935v1 (30 October 2006).

\bibitem[dFR82]{dFR82}
H.~de~Fraysseix and P.~Rosenstiehl,
A depth-first-search characterization of planarity,
\emph{Ann. Discrete Math.} \textbf{13} (1982), 75--80.

\bibitem[dFR85]{dFR85}
H.~de~Fraysseix and P.~Rosenstiehl,
A characterization of planar graphs by Tr\'emaux orders,
\emph{Combinatorica} \textbf{5} (1985), 127--135.

\bibitem[GM13]{GM13}
N.~Robertson and P.~D.~Seymour,
Graph Minors. XIII. The disjoint paths problem,
\emph{J. Combin. Theory Ser. B} \textbf{63} (1995), 65--110.

\bibitem[HJ26]{HJ26}
M.~Hajiaghayi and M.~JafariRaviz,
Extracting a crossing in time $O(nm)$,
note, November 2026, file \path{docs/crossing/crossing_extractor.tex} of the
repository \texttt{GraphMinor}.

\bibitem[HSS08]{HSS08}
A.~A.~Hagberg, D.~A.~Schult and P.~J.~Swart,
Exploring network structure, dynamics, and function using NetworkX,
in \emph{Proceedings of the 7th Python in Science Conference (SciPy 2008)},
2008, 11--15.

\bibitem[HT74]{HT74}
J.~Hopcroft and R.~Tarjan,
Efficient planarity testing,
\emph{J. ACM} \textbf{21} (1974), 549--568.

\bibitem[Sh80]{Sh80}
Y.~Shiloach,
A polynomial solution to the undirected two paths problem,
\emph{J. ACM} \textbf{27} (1980), 445--456.

\end{thebibliography}}

\end{document}
